
\def\PUDcodigo{ACE009}

\if\COMPILAR0
   \def\PUDcodacad{04507.17}
\else
   \def\PUDcodacad{04506.17}
\fi

\def\PUDdisciplina{HST}

\def\PUDsemestre{S4}

\def\PUDhoras{40}

%NOVOS ---------------------------------------------------
\def\PUDhorasTeoria{40}
\def\PUDhorasPratica{0}

\def\PUDinterECA{---}

\def\PUDatuaisECA{---}

\def\PUDrecursos{Projetor multimídia; Quadro branco e pincel.}

%----------------------------------------------------------

\def\PUDcreditos{2}

\def\PUDprerequisito{---}

\def\PUDpreacad{---}

\def\PUDobjetivos{Compreender os princípios fundamentais de segurança e higiene no trabalho na engenharia.  Tomar conhecimento das condições de higiene e segurança no trabalho, como também a proteção coletiva e individual e proteção contra incêndio.  Participar de programas de prevenção de riscos.  Conhecer a legislação especifica e normas técnicas.}

\def\PUDementa{Legislação e normas.  Implantação da segurança do trabalho.  Controle estatístico de acidentes.  Equipamentos de proteção individual e coletivo.  Iluminação.  Ruído.  Calor.  Frio.  Umidade.  Sinalização e cor.  Condições sanitárias e de confronto.}

\def\PUDconteudo{ &  Legislação sobre segurança e saúde no trabalho: A Constituição;  A Consolidação das Leis Trabalhistas (CLT);  As Normas Regulamentadoras (NRs).  
\\ 
 &  A Comissão Interna de Prevenção de Acidentes – CIPA: Atribuições;  Organização e funcionamento.  
\\ 
 &  O serviço de Engenharia e Medicina do Trabalho: A formação;  Atribuições do médico do trabalho;  Atribuições do Engenheiro de Segurança.  
\\ 
 &  Local de trabalho: Riscos graves e interdição;  Inspeção;  Investigação das causas dos acidentes;  As causas dos acidentes;  Ato inseguro e condição insegura;  Proteção de Máquinas e Equipamentos;  Dispositivos de acionamento e parada;  Riscos com eletricidade.  
\\ 
 &  Esforço físico e movimentação de materiais: O esforço físico e as lesões;  Cuidados e orientações preventivas;  Normas legais/ergonomia;  Consequências do excesso de trabalho;  Duração da jornada de trabalho e ritmo de trabalho;  Período de repouso;  
\\ 
 &  Proteção contra incêndios: Como evitar o fogo / Como combater o incêndio;  Classes de fogo e métodos de extinção;  Saídas de emergência.  
\\ 
 &  Insalubridade e riscos no trabalho: Avaliação dos limites de tolerância;  Ruído, calor, iluminação, riscos químicos;  O adicional de insalubridade.                                                                                                      \\ 
 &  Medidas de proteção contra riscos ocupacionais: Proteção individual e proteção coletiva;  Implantação de um EPI;  Normas Legais sobre EPI.  
\\ 
 &  Periculosidade: Explosivos;  Inflamáveis;  Eletricidade;  Radioatividade.  
\\ 
 &  Noções de primeiros socorros. }

\def\PUDmetodo{Aulas expositórias que podem ser teóricas e/ou práticas, onde as práticas no laboratório serão marcadas no decorrer da disciplina.}

\def\PUDavalia{A avaliação será feita com aplicação de provas teóricas e/ou práticas, além da possibilidade de inclusão de trabalhos e seminários no decorrer da disciplina.}

\def\PUDbibbasica{ & \href{www.biblioteca.ifce.edu.br/index.asp?codigo_sophia=61973}{Oliveira}, Cláudio Antonio Dias de.  Manual Prático de Saúde e Segurança do Trabalho.  2º ed.  São Caetano do Sul, SP: Yendis, 2012.  433p.  ISBN: 9788577282593.
%& ARAÚJO, Giovanni Moraes de.  Legislação de segurança e saúde ocupacional: Normas regulamentadoras do Ministério do Trabalho e Emprego.  Rio de Janeiro: GVC Gerenciamento Verde Consultoria, 2006.  %ISBN 9788599331026
\\ 
& \href{www.biblioteca.ifce.edu.br/index.asp?codigo_sophia=42081}{Couto}, Hudson de Araújo.  Gerenciando a LER e os DORT nos tempos atuais.  Belo Horizonte, MG: Ergo, 2007.  492p.  ISBN: 978859959028.
%\\ 
 %&   PACHECO JUNIOR, Waldemar.  Gestão da segurança e higiene do trabalho: Contexto estratégico, análise ambiental, controle e avaliação das estratégias.  São Paulo: Ed.  Atlas, 2000.  ISBN %9788522424368
\\ 
 &   \href{www.biblioteca.ifce.edu.br/index.asp?codigo_sophia=39900}{Seiffert}, Mari Elizabete Bernardini.  Sistemas de gestão ambiental (ISO 14001) e saúde e segurança ocupacional (OHSAS 18001): vantagens da implantação integrada.  2ª ed.  São Paulo, SP: Atlas, 2010.  201p.  ISBN: 9788522451111. }
%\\ 
 %&   ARAÚJO,  W.  T.  Manual de Segurança do Trabalho.  1ª ed. , Editora: DCL, 2010.

\def\PUDbibcomplementar{ &  \href{www.biblioteca.ifce.edu.br/index.asp?codigo_sophia=40795}{Peraire}, José M. Paris.  Manual do montador de quadros elétricos: características dos materiais, sua qualidade, sua forma de construção.  São Paulo, SP: Hemus, 2004.  233p.  ISBN: 8528904040.
 %& Saliba, Tuffi Messias.  Curso básico de segurança e higiene ocupacional.  São Paulo, SP: LTr, 2004.  452p.  ISBN: 853610516.
 %& Iida, Itiro.  Ergonomia: projeto e produção.  2º ed.  São Paulo, SP: Blucher, 2005.  614p.  ISBN: 978852120354.
\\
 &  \href{www.biblioteca.ifce.edu.br/index.asp?codigo_sophia=40790}{Barros}, B. F. de; Borelli, R.; Guimarães, E. C. De A.  NR-10: guia prático de análise e aplicação. 3º ed.  São Paulo, SP: Érica, 2014.  204p.  ISBN: 9788536502748.
\\
 &  \href{www.biblioteca.ifce.edu.br/index.asp?codigo_sophia=24014}{Pepplow}, Luiz Amilton.  Segurança do trabalho.  Curitiba, PR: Base Editorial, 2010.  256p.  ISBN: 9788579055430.
\\
 &  \href{www.biblioteca.ifce.edu.br/index.asp?codigo_sophia=4676}{Zocchio}, Alvaro.  Prática da prevenção de acidentes: o ABC da segurança do trabalho.  E-book.  São Paulo, SP: Atlas, 1965.  227p.  ISBN: 9788522472994.
\\
 &  \href{www.biblioteca.ifce.edu.br/index.asp?codigo_sophia=58184}{Rossete}, Celso Augusto.  Segurança e higiene do trabalho.  E-book.  Pearson.  186p.  ISBN: 9788543012216. }
 %& Travassos, Geraldo.  Guia prático de medicina do trabalho: a influência dos ambientes de trabalho na saúde de nossos pacientes: o que todo médico precisa saber sobre saúde ocupacional.  São Paulo, SP: LTr, 2003.  110p.  ISBN: 8536104783.
%& CARUSO, Marina.  Um perigo real.  In: Isto é, nº1686.  São Paulo.  Ed.  Três, 23 de janeiro de 2002. 
%\\ 
 %&  MAENO, Mara et al.  Lesões por Esforços Repetitivos (LER) e distúrbios Osteomusculares.  Brasília: Ministério da saúde, 2001.

\def\PUDautor{Fabrício Bandeira da Silva}

\def\PUDdata{2014-05-19}

\def\PUDrevisao{2}

\def\PUDrevisor{Samuel Vieira Dias}

\def\PUDdatarevisao{2017-05-11}

\PUDcorpo
